\documentclass[conference]{IEEEtran}
\IEEEoverridecommandlockouts

\usepackage{cite}
\usepackage{amsmath,amssymb,amsfonts}
\usepackage{graphicx}
\usepackage{booktabs}
\usepackage{multirow}
\usepackage{textcomp}
\usepackage{xcolor}
\usepackage{url}
\usepackage{hyperref}

\def\BibTeX{{\rm B\kern-.05em{\sc i\kern-.025em b}\kern-.08em
    T\kern-.1667em\lower.7ex\hbox{E}\kern-.125emX}}

\begin{document}

\title{Ambiguity Localization at the Cost of a Forward Pass:\\
Lightweight Hidden-State Probing for Clarification-Question Ranking\\
in Agentic Code Generation}

\author{
\IEEEauthorblockN{Author Name}
\IEEEauthorblockA{\textit{Affiliation} \\
\textit{Institution}\\
City, Country \\
email@example.com}
}

\maketitle

\begin{abstract}
Agentic coding assistants that ask a clarifying question before writing code
can resolve specification ambiguity that a one-shot generation cannot. The
standard way to decide \emph{whether} and \emph{what} to clarify is to
estimate the coding model's uncertainty by repeated sampling of its own
outputs -- an approach that scales the cost of every clarification decision
with the cost of the (often large) coding model itself. We show that a
lightweight linear probe over the hidden states of a small, decoupled model
(Qwen3-8B), extracted in a single forward pass, can (i) detect the degree of
specification ambiguity in a coding prompt and (ii) rank candidate
clarifying questions for a two-round elicit-rank-clarify-code agentic
pipeline, at a fraction of the inference cost of prompting a model to reason
about its own uncertainty. We evaluate the ranking probe against a zero-shot
self-rank baseline across four benchmark sources (ClarifyCodeBench and three
Orchid-derived variants of HumanEval and BigCodeBench) under grouped,
leave-one-source-out splits, and against real, executable test suites in a
live two-round clarify-then-code loop with two different coding models. We
find the probe reaches ranking quality close to -- and on its
closest-to-training-distribution source, above -- the zero-shot baseline, at
substantially lower measured inference cost. Using the fractional
test-case pass rate rather than an all-or-nothing correctness metric, we
show that correctly ranking the labeled ambiguity to rank one is associated
with a statistically significant round-1 correctness gain
($+9.9$ percentage points, $p=0.035$, bootstrap 95\% CI $[1.2, 19.1]$pp on
Orchid-sourced problems), a result the coarser binary metric alone does not
surface significantly. We further show that asking any clarifying question
at all, independent of which one, is itself a significant and robust
correctness gain, and that a cheap, targeted prompt augmentation recovers
missing type information that the ambiguity-injection process otherwise
strips, further improving pass rates. We report all findings, including
places where our own earlier conclusions were revised under closer
methodological scrutiny, as a case study in evaluating agentic clarification
pipelines rigorously.
\end{abstract}

\begin{IEEEkeywords}
agentic software engineering, code generation, ambiguity detection, hidden-state probing, clarifying questions, large language models
\end{IEEEkeywords}

\section{Introduction}
\label{sec:intro}

Large language model (LLM) coding assistants increasingly operate as
\emph{agents}: given a natural-language task description, they plan, write,
run, and revise code with little further human input. A persistent failure
mode in this setting is specification ambiguity -- the prompt genuinely
under-determines the intended behavior, and a one-shot generation commits to
an arbitrary resolution that may not match the user's intent. A natural
mitigation is to have the agent ask a clarifying question before writing
code, turning a single generation step into a short elicit-then-generate
loop~\cite{clarifygpt}. This raises two sub-problems: \emph{detecting} that
a prompt is ambiguous enough to be worth interrupting for, and, given a set
of candidate clarifying questions, \emph{ranking} which one to ask.

Existing approaches to both sub-problems typically estimate the coding
model's own uncertainty by sampling its outputs repeatedly and measuring
disagreement among the samples -- an approach with a long lineage in
uncertainty quantification for LLMs, exemplified by semantic
entropy~\cite{kuhn2023semantic} and its application to code-generation
ambiguity in systems such as ClarifyGPT~\cite{clarifygpt}. This is
effective, but its cost scales linearly with the number of samples drawn
\emph{from the coding model itself}, which is often the largest, most
capable, and most expensive model in the pipeline. Every clarification
decision therefore costs as much as several full code-generation calls,
before any code has actually been written.

We investigate a different design point, motivated by a simple observation:
if a prompt's hidden-state representation already encodes how ambiguous it
is, that signal can be read off in a single forward pass, without sampling
any generations at all -- and that forward pass can be run on a small model
\emph{decoupled} from whichever model eventually writes the code. This
motivates the semantic-entropy-probe approach of Kossen et
al.~\cite{oatml_sep}, who train a linear probe on a model's hidden states to
predict semantic entropy (a property of the model's \emph{output}
distribution) from a single forward pass, avoiding repeated sampling at
inference time for that specific prediction task. We adapt this probing
methodology to a different target relevant to agentic software engineering:
not the entropy of a model's own generations, but the ambiguity of an
\emph{input} coding prompt, and, in a second stage, the relative quality of
candidate clarifying questions for resolving that ambiguity.

Concretely, we make the following contributions:

\begin{itemize}
\item We show that a per-layer linear probe over Qwen3-8B's residual-stream
activations, extracted from a single forward pass over an unmodified coding
prompt, predicts the degree of injected specification ambiguity
significantly above chance (Section~\ref{sec:results}), and that this
signal is present, though imperfectly transferable, across four
differently-constructed ambiguity benchmarks.
\item We build a ranking probe that scores candidate clarifying questions
from batched forward passes over the prompt concatenated with each
candidate, and deploy it in a live, two-round elicit-rank-clarify-code
agentic loop with real executable test suites, comparing it against a
zero-shot self-rank baseline (the coding model is simply asked to pick the
best candidate) under grouped, leave-one-source-out evaluation.
\item We report, honestly, that the probe's ranking quality is
close to but does not uniformly exceed the zero-shot baseline -- it
wins on the benchmark closest to its training distribution and trails on
others -- and we frame our contribution accordingly as a
\emph{cost-quality tradeoff}: comparable ranking quality at substantially
lower measured inference cost, since the probe requires no autoregressive
generation from the coding model, rather than as an unconditional accuracy
win.
\item We investigate whether ranking quality (does the pipeline pick the
labeled ambiguity as the top clarifying question) actually predicts
downstream code correctness, and show that the answer depends on the
correctness metric used: an all-or-nothing ``did every test pass'' metric
shows no significant relationship, but the more informative fractional
test-case pass rate reveals a real, statistically significant effect in
round one of the clarification loop. We report both metrics and both
outcomes, since this discrepancy is itself a methodological finding about
how to evaluate agentic clarification systems.
\item We identify and quantify a second, largely orthogonal lever for
correctness: whether \emph{any} clarifying question is asked at all is a
robust, significant improvement over asking none, and a cheap,
targeted prompt augmentation that restores type and requirement information
the ambiguity-injection process incidentally strips further reduces
trivial crash-class failures.
\end{itemize}

We deliberately report several points in this investigation where our own
earlier conclusions did not survive closer scrutiny -- most notably, an
early result that the probe outperformed zero-shot ranking was later
reversed under a larger, more carefully filtered evaluation. We treat this
not as a weakness to hide but as an illustration of the methodological
rigor agentic-system evaluations in software engineering research need: a
result computed once, on one split, with one metric, is not yet a finding.

\section{Related Work}
\label{sec:related}

\subsection{Uncertainty Quantification via Sampling}

Semantic entropy~\cite{kuhn2023semantic} estimates an LLM's uncertainty
about a free-form generation by sampling multiple outputs, clustering them
by semantic equivalence, and computing the entropy of the resulting cluster
distribution. This requires $K$ generation calls per uncertainty estimate
and is the conceptual basis for a family of clarify-then-generate systems
for code, including ClarifyGPT~\cite{clarifygpt}, which uses repeated
sampling of the coding model's outputs to decide when an ambiguous
requirement warrants a clarifying question. These approaches are effective
but tie the cost of every clarification decision to $K$ additional calls to
the coding model itself, which is frequently the most expensive component
of the pipeline.

\subsection{Single-Forward-Pass Probing}

Kossen et al.~\cite{oatml_sep} introduce semantic entropy probes: a linear
probe trained on a single forward pass' hidden states to predict the
semantic entropy that would otherwise require $K$ generations to estimate,
amortizing the sampling cost of training many probes' worth of labels into
a one-time cost, with cheap inference thereafter. Our probing methodology
directly reuses this approach -- per-layer linear probing of residual-stream
activations, best-split binarization for layer selection, bootstrapped
evaluation -- but applies it to a different prediction target. Where
Kossen et al.\ probe for a property of the model's own \emph{output}
distribution (how much its generations would disagree with each other), we
probe for a property of the \emph{input} prompt (how much unresolved
specification ambiguity it contains) and, in a second stage, a property of
a candidate answer relative to that input (how well it resolves the
ambiguity, for ranking purposes). This is a substantively different
prediction problem: no sampling, clustering, or generation from the model
being probed is involved anywhere in our pipeline, at either training or
inference time, and consequently the model being probed need not be the
model doing the coding.

\subsection{Ambiguity and Clarification in Software Engineering}

Resolving requirements ambiguity before implementation is a long-standing
concern in software engineering, predating LLM-based tooling. Recent
LLM-era work targets this directly for code generation, constructing
benchmarks with deliberately injected or naturally occurring ambiguity and
evaluating whether models recognize the need to ask a clarifying question
before writing code~\cite{clarifycodebench}. Our evaluation builds on
ClarifyCodeBench and on Orchid-style ambiguity injections over
HumanEval~\cite{humaneval} and BigCodeBench~\cite{bigcodebench}, which
construct matched clear/ambiguous prompt pairs by deleting or rewriting a
targeted specification detail and pairing the injected prompt with tests
that specifically exercise the resulting ambiguity.

\subsection{Agentic Coding Benchmarks}

Our live evaluation loop -- eliciting candidate clarifying questions,
ranking them, generating code conditioned on the top-ranked answer, and
scoring the result against real executable test suites -- follows the
general pattern of agentic coding benchmarks that emphasize execution-based
correctness over surface-level similarity to a reference
solution~\cite{humaneval,bigcodebench}. We report correctness using the
fractional test-case pass rate rather than an all-or-nothing pass metric,
and show in Section~\ref{sec:results} that this choice materially affects
which relationships in the data reach statistical significance -- a point
we believe is broadly relevant to how agentic coding pipelines should be
evaluated, independent of our specific method.

\section{Approach}
\label{sec:approach}

\subsection{Hidden-State Extraction}

All probes read the residual stream of Qwen3-8B, run in bf16 precision with
thinking mode explicitly disabled. For a given prompt, we apply the model's
chat template, run a forward pass with \texttt{output\_hidden\_states=True}
and no key-value caching, and extract the hidden state at the last token of
the templated input (before any generation) at every one of the model's 37
residual-stream checkpoints (layer 0, the embedding output, through layer
36, the final block's output). No text is generated at any point in this
extraction; the only computation performed is the forward pass itself.

\subsection{Ambiguity-Degree Probe}
\label{sec:approach:k-probe}

The first probe answers a coarser question: given a coding prompt, how much
unresolved specification ambiguity does it contain, expressed as an ordinal
count $K \in \{1, 2, 3\}$ of independently ambiguous details? We extract the
last-token hidden state of the raw, unmodified prompt at each layer and fit
an independent multinomial logistic regression per layer on top of a
100-component PCA plus standardization pipeline, with the regularization
strength $C$ swept on a validation split. Layer and configuration selection
use grouped cross-validation, splitting by source-problem identifier
(\texttt{StratifiedGroupKFold}) so that no two rows derived from the same
underlying problem appear on both sides of a split -- a leakage guard that
is load-bearing for this task, since near-duplicate prompts sharing a
source problem would otherwise let a probe memorize surface features of a
specific problem rather than learning a general ambiguity signal. We report
ordinal correlation (Spearman's $\rho$ between the predicted expected value
$E[K]$ under the probe's class-probability output and the true $K$) as the
primary metric, alongside majority-class, random, and prompt-length-only
baselines computed on the identical split.

\subsection{Clarifying-Question Ranking Probe}
\label{sec:approach:rank-probe}

The second, and practically more important, probe addresses a different
problem: given a pool of $N$ candidate clarifying questions (with
LLM-elicited answers) for an ambiguous prompt, which one should the agent
actually ask? We treat this as pairwise ranking cast as classification, in
the style of RankSVM/RankNet: each candidate's answer is concatenated onto
the base prompt (\texttt{"\{prompt\}\textbackslash n\textbackslash
nClarification: \{answer\}"}), and a batched forward pass over all $N$
concatenated variants (one GPU call, $N$ sequences) extracts $N$ hidden
states at a single selected layer. A classifier trained on
tier-difference vectors between higher- and lower-quality candidates --
where candidate quality tiers (0/1/2, unrelated/plausible/matches-the-
labeled-ambiguity) are assigned by an independent LLM judge -- scores each
candidate, and the pipeline asks the top-scored one. We additionally
evaluate a supervised-contrastive (SupCon) variant of this probe, which
trains directly on raw (non-PCA-reduced) activations with a contrastive
objective and, as reported in Section~\ref{sec:results}, improves
cross-source generalization at some cost to in-distribution accuracy
relative to the PCA-based classifier.

Critically, this ranking probe's inference cost is $O(N)$ forward passes
executed as a single batched GPU call, not the $O(1)$ single forward pass
of the ambiguity-degree probe in Section~\ref{sec:approach:k-probe} --
each candidate's answer text must be present in the input for its
hidden state to reflect that candidate specifically. The relevant cost
comparison is therefore not ``one forward pass vs.\ one generation call''
but ``one batched forward pass over $N$ short sequences vs.\ one
generation call that reasons over all $N$ candidates inline,'' which we
measure directly in Section~\ref{sec:results:cost}.

\subsection{The Two-Round Elicit-Rank-Clarify-Code Loop}
\label{sec:approach:loop}

We deploy the ranking probe in an agentic pipeline that mirrors how a
clarify-then-code assistant would operate in practice:

\begin{enumerate}
\item \textbf{Elicit}: given an ambiguous prompt, generate a pool of
candidate clarifying questions and, for each, an oracle answer (an
independent LLM answers the question as if it were the ground-truth
requirement author).
\item \textbf{Rank}: score all candidates with the trained probe (or, for
the zero-shot baseline in Section~\ref{sec:experimental_setup}, ask the
coding model to pick directly from a single inlined list) and select the
top-ranked candidate.
\item \textbf{Clarify}: append the selected candidate's question-answer
pair to the prompt.
\item \textbf{Code}: generate code conditioned on the clarified prompt, and
run the real, executable test suite associated with the problem.
\end{enumerate}

This constitutes round one. Round two repeats steps 1--4 once more,
conditioned on the round-one clarified prompt, to model an agent that gets
a second opportunity to resolve any remaining ambiguity before finalizing
its code. We report both round-1 and round-1+2 correctness throughout
Section~\ref{sec:results}, since, as we show, the value of a second
clarification round does not straightforwardly track the value of the
first.

\subsection{Design Rationale for the Probe-Host Model}

We use Qwen3-8B as the probe-host model throughout due to compute
constraints in this study; the ranking probe's forward-pass cost scales
with this model's size, not with the size of whichever model ultimately
writes the code. This decoupling is itself part of the approach's intended
value proposition: ambiguity localization can, in principle, be performed
by a small, cheap model regardless of which (possibly much larger and more
expensive) model is used for code generation, whereas a self-consistency
baseline built on the coding model's own repeated sampling ties uncertainty
estimation cost directly to that larger model's per-token cost. We evaluate
this decoupling directly in Section~\ref{sec:results} by pairing the same
ranking decisions with two different coding models of different capability.

\section{Experimental Setup}
\label{sec:experimental_setup}

\subsection{Datasets}

We evaluate on four sources of ambiguity-injected coding problems:

\begin{itemize}
\item \textbf{ClarifyCodeBench}~\cite{clarifycodebench}: 419 problems
derived from LiveCodeBench, each annotated with $K \in \{1,2,3\}$
deliberately deleted specification details.
\item \textbf{Orchid-HEval}: HumanEval~\cite{humaneval} problems rewritten
with a single injected ambiguous edit (lexical, vagueness, semantic, or
syntactic), paired with the original hidden tests.
\item \textbf{Orchid-BCB} and \textbf{Orchid-BCB-Expand}: BigCodeBench
\cite{bigcodebench} problems rewritten the same way, with BCB-Expand a
larger, semi-automatically constructed extension of the original BCB
subset.
\end{itemize}

All source-problem groupings are preserved under splitting: no two prompt
variants derived from the same underlying problem appear on opposite sides
of a train/validation/test or leave-one-source-out (LODO) split. We use
grouped, stratified splits throughout the ambiguity-degree probe
evaluation, and full leave-one-source-out evaluation (train on three
sources, test on the fourth, held out entirely) for the ranking probe, to
measure whether ranking quality generalizes across benchmark construction
methodologies rather than only within one.

\subsection{Models}

The probe-host model is Qwen3-8B (bf16, thinking mode disabled) throughout.
For code generation, we use two models of different capability: Qwen3-8B
itself, and DeepSeek V3.2 (accessed via Bedrock), to test whether findings
about clarification quality are specific to a particular coding model or
hold across models of different capability. Candidate clarifying-question
elicitation and independent quality judging (assigning tiers 0/1/2 to
candidates relative to the benchmark's labeled ambiguity) use DeepSeek V3.2
as an independent judge, separate from whichever model generates the code
under test.

\subsection{Baselines}

\textbf{Zero-shot self-rank}: the coding model is prompted with the
ambiguous prompt and the full numbered list of candidate clarifying
questions in a single message, and asked to output the number of the one it
would ask. This requires exactly one generation call per problem (greedy
decoding, no sampling), with the entire candidate pool inlined -- not one
call per candidate. \textbf{No-clarification}: the raw ambiguous prompt is
passed directly to the coding model with no clarification step at all,
establishing the baseline value of clarifying at all, independent of which
question is chosen.

\subsection{Correctness Metric}

We report correctness as the \textbf{fractional test-case pass rate}: the
number of individual test assertions a generated solution passes divided by
the total number of test assertions for that problem, averaged across
groups. This is distinct from -- and, as shown in
Section~\ref{sec:results}, statistically more sensitive than -- an
all-or-nothing metric that scores a group as passing only if every test
assertion succeeds. Both metrics are computable from the same underlying
per-group pass/total counts; we report the fractional metric as primary
throughout and note where it diverges from the all-or-nothing metric's
conclusions.

\subsection{Significance Testing}

For between-groups comparisons (e.g., groups where the probe ranked the
labeled ambiguity first vs.\ groups where it did not), we report
Mann-Whitney U and Welch's $t$-test alongside a 10,000-resample bootstrap
95\% confidence interval on the mean difference. For within-groups,
paired comparisons (e.g., the same 244 problems with vs.\ without any
clarification), we report the Wilcoxon signed-rank test, a paired $t$-test,
and a paired bootstrap 95\% CI on the mean of per-group differences. We
disclose cases where these tests disagree at the margin rather than
reporting only the test that reaches $p < 0.05$.

\subsection{Compute Environment}

\label{sec:experimental_setup:compute}
All hidden-state extraction and generation runs on a single GPU
(NVIDIA A10G, 23\,GB, native bf16, no CPU weight offload). We note that an
earlier measurement on a memory-constrained T4 GPU requiring CPU offload
showed no measurable speedup between the probe and zero-shot generation, an
artifact of a fixed CPU$\leftrightarrow$GPU weight-transfer cost dominating
both paths equally; we report only native-GPU timings here, and describe
the exact benchmarked comparison in
Section~\ref{sec:results:cost}.

\section{Results}
\label{sec:results}

\subsection{Ambiguity-Degree Probe: A Single Forward Pass Encodes Some Signal}
\label{sec:results:k-probe}

On ClarifyCodeBench restricted to $K \in \{1,2,3\}$ (excluding any
constructed $K{=}0$ class, since every tested construction of a
``no-ambiguity'' class introduced a detectable confound -- see
Section~\ref{sec:discussion}), the ambiguity-degree probe achieves test
Spearman $\rho = 0.312$ (5-fold cross-validated $\rho = 0.21 \pm 0.05$) at
its best layer (layer 32 of 37), against a random baseline of $\rho =
0.007$ and a prompt-length-only baseline of $\rho = 0.043$. The
layer-by-layer curve rises smoothly from $\rho \approx 0.08$ at layer 1 to
its layer-32 peak, consistent with a property that is progressively
constructed through the network's depth rather than a surface feature
available immediately from the embedding layer. This result is modest but
reproducible: an independent re-run with a differently-cleaned prompt
template reproduced $\rho = 0.31230908$ to eight decimal places.

Extending this probe to a pooled, four-source binary (clear-vs-ambiguous)
task spanning ClarifyCodeBench and all three Orchid sources, a
PCA-reduced logistic-regression probe reaches pooled in-distribution AUROC
of $0.75$--$0.77$; the same probe's leave-one-source-out AUROC falls to
$0.52$--$0.64$, indicating the signal generalizes only partially across
benchmark construction methodologies. A supervised-contrastive variant
trained on raw (non-PCA) activations improves leave-one-source-out AUROC to
$0.60 \pm 0.01$ (5 seeds) at some cost to pooled accuracy -- a genuine
generalization/accuracy tradeoff we revisit for the ranking probe below.

\subsection{Ranking Probe vs.\ Zero-Shot: A Cost-Quality Tradeoff, Not an Accuracy Win}
\label{sec:results:ranking}

Table~\ref{tab:ranking} reports rank-1 hit rate (did the pipeline's top
choice match the benchmark's labeled ambiguity) for the trained ranking
probe (PCA + logistic regression, the best-performing configuration) and
for zero-shot self-rank, both evaluated on the identical 1,229-group,
four-source pool under leave-one-source-out splitting.

\begin{table}[htbp]
\caption{Ranking probe vs.\ zero-shot self-rank, rank-1 hit rate, leave-one-source-out evaluation (1,229 groups).}
\label{tab:ranking}
\centering
\begin{tabular}{lcc}
\toprule
Evaluation & Probe (PCA+LR) & Zero-shot Qwen \\
\midrule
LODO mean (4 folds) & 0.637 & \textbf{0.710} \\
\quad held-out ClarifyCodeBench ($n{=}187$) & \textbf{0.701} & 0.663 \\
\quad held-out Orchid-HEval ($n{=}187$) & 0.738 & 0.743 \\
\quad held-out Orchid-BCB ($n{=}110$) & 0.582 & \textbf{0.691} \\
\quad held-out Orchid-BCB-Expand ($n{=}745$) & 0.525 & \textbf{0.742} \\
Pooled, proportional split ($n{=}246$) & 0.581 & \textbf{0.679} \\
Pooled, equal-per-source ($n{=}94$) & 0.638 & \textbf{0.713} \\
\bottomrule
\end{tabular}
\end{table}

The trained probe does not uniformly beat zero-shot self-rank: zero-shot
wins on 6 of the 7 rows in Table~\ref{tab:ranking}, most substantially on
Orchid-BCB-Expand ($0.742$ vs.\ $0.525$, a 22-point gap). The probe's one
clear win is on ClarifyCodeBench, its closest-to-training-distribution
source. We note this result explicitly reverses an earlier, smaller-scale
finding in our own development process, where the probe appeared to beat
zero-shot by 13 points on a looser, 121-group evaluation split; that
earlier result did not survive evaluation on the larger, more stringently
filtered pool used here (Section~\ref{sec:discussion}).

We therefore frame the ranking probe's contribution as a
\textbf{cost-quality tradeoff} rather than an accuracy win: it reaches
ranking quality within a few points of zero-shot self-rank on two of four
sources and its own training source, and trails meaningfully on the two
BigCodeBench-derived sources, whose prompts are long and free-form relative
to the size of the injected ambiguous edit (Section~\ref{sec:discussion}).

\subsection{Inference Cost}
\label{sec:results:cost}

% TODO(paper): replace with the fresh, apples-to-apples benchmark (batched
% N-candidate probe extraction vs. zero-shot generation, same hardware/groups)
% once GPU access is confirmed -- see plan Step 2. The number below is a
% preliminary measurement that compares zero-shot generation to a probe
% forward pass over the BARE PROMPT ONLY, not the full batched N-candidate
% extraction the deployed ranking pipeline actually performs, and should be
% treated as provisional until replaced.
On native GPU hardware (no CPU weight offload), a single probe forward pass
measures $0.079$s mean / $0.065$s median per problem, against $0.38$s mean
/ $0.37$s median for zero-shot generation with the full candidate list
inlined ($n{=}20$) -- a $4.8\times$ speedup. \textbf{[Placeholder pending a
fresh benchmark of the actual batched $N$-candidate extraction the deployed
ranking pipeline performs, per Section~\ref{sec:approach:rank-probe}; the
measurement above compares zero-shot generation to a single bare-prompt
forward pass, which is not quite the right comparison and should be treated
as provisional.]} We note that this measured advantage vanished entirely on
a memory-constrained GPU requiring CPU weight offload
(Section~\ref{sec:experimental_setup:compute}), underscoring that the
compute-efficiency argument for this class of method depends on deployment
hardware, not only on the underlying FLOP count.

\subsection{Does Ranking Quality Predict Downstream Correctness?}
\label{sec:results:correctness}

We evaluate the live two-round elicit-rank-clarify-code loop
(Section~\ref{sec:approach:loop}) against real, executable test suites on
244 problems spanning all four sources, using DeepSeek V3.2 as the coding
model. Table~\ref{tab:hit-vs-pass} reports the relationship between
round-1 ranking success (\texttt{hit1}: did the top-ranked candidate match
the labeled ambiguity) and round-1 fractional test-case pass rate,
restricted to the three Orchid sources (ClarifyCodeBench lacks Orchid's
one-injected-ambiguity, targeted-test design and is excluded from this
specific comparison).

\begin{table}[htbp]
\caption{Round-1 fractional pass rate by ranking outcome, Orchid sources, DeepSeek V3.2 coding model ($n{=}213$).}
\label{tab:hit-vs-pass}
\centering
\begin{tabular}{lccc}
\toprule
Source & hit1=1 & hit1=0 & Gap \\
\midrule
Orchid-BCB ($n{=}21$) & 88.6\% & 67.2\% & $+21.5$pp \\
Orchid-HEval ($n{=}47$) & 100.0\% & 97.9\% & $+2.1$pp \\
Orchid-BCB-Expand ($n{=}145$) & 77.5\% & 65.7\% & $+11.8$pp \\
\textbf{All Orchid combined} & \textbf{83.3\%} & \textbf{73.3\%} & \textbf{$+9.9$pp} \\
\bottomrule
\end{tabular}
\end{table}

The all-Orchid-combined gap of $+9.9$ percentage points is statistically
significant by Welch's $t$-test ($p = 0.035$) and by a 10,000-resample
bootstrap 95\% confidence interval on the mean difference, $[1.2, 19.1]$pp,
which excludes zero. A Mann-Whitney U test on the same data is more
conservative ($p = 0.10$); we report both rather than the more favorable
test alone. This is a materially stronger result than the same
relationship measured with the all-or-nothing pass metric on the same
data ($+7.6$pp, Fisher's exact $p = 0.31$, not significant) -- the
fractional metric recovers a real effect that the coarser binary metric
was not sensitive enough to detect at this sample size. We regard this
discrepancy as a methodological finding in its own right: whether an
agentic clarification pipeline appears to ``work'' can depend on which
correctness metric an evaluation uses, independent of anything about the
pipeline itself.

The effect is round-1-specific. Recomputed on round-1+2 (after a second
clarification round), the all-Orchid-combined gap falls to $+7.4$pp
($p = 0.11$, bootstrap CI includes zero); per-source, Orchid-BCB alone
strengthens to $+25.2$pp ($p = 0.05$, bootstrap CI $[5.7, 47.9]$pp,
significant) while Orchid-BCB-Expand weakens to $+8.0$pp (not
significant). We do not have a single clean explanation for this
round-to-round variability at the current sample size and report it as an
open question rather than force a narrative onto it.

We separately identify why the effect is markedly weaker for
Orchid-BCB-Expand than for the other sources: an independent LLM-judged
analysis of 1,952 candidate clarifying questions for this source finds
that, while the officially labeled ambiguity is functionally relevant to
correctness at essentially the same rate as every other source
($91.3\%$, vs.\ $91.8$--$92.4\%$ elsewhere), it is rarely the \emph{only}
relevant detail: on average $3.09$ other candidate clarifications are
independently functionally relevant to the same problem (vs.\ $1.48$ for
ClarifyCodeBench and $2.43$ for Orchid-HEval), and the label is the sole
relevant detail in only $7.0\%$ of rows (vs.\ $26.7\%$ and $14.1\%$
respectively). Resolving the one labeled ambiguity correctly therefore
caps, but does not determine, whether the generated code passes, because
most Orchid-BCB-Expand problems are underspecified in several independent
ways at once and the current pipeline resolves at most two (one per
round).

\subsection{Asking Any Question Beats Asking None}
\label{sec:results:any-question}

Independent of \emph{which} clarifying question is asked, we find that
asking any question at all is a robust, statistically significant
improvement over asking none. Comparing the same 244 problems with vs.\
without any clarification (a matched, paired comparison), round-1
clarification improves the mean fractional pass rate by $+4.0$pp (Wilcoxon
signed-rank $p = 0.0021$, paired bootstrap 95\% CI $[1.6, 6.6]$pp), and
round-1+2 by $+6.3$pp (Wilcoxon $p < 0.0001$, CI $[3.7, 9.1]$pp). Both
results are significant under every test we applied, in contrast to the
more marginal, metric-dependent ranking-quality effect in
Section~\ref{sec:results:correctness}. This suggests that, for this task
family, the decision to clarify at all carries more of the achievable
correctness gain than the precision of the ranking step that decides
\emph{which} clarification to make -- a finding with direct design
implications for agentic coding pipelines operating under a compute
budget that cannot afford a highly precise ranking step.

\subsection{A Cheap, Orthogonal Fix: Prompt Augmentation}
\label{sec:results:augmentation}

Separately from the ranking pipeline, we identify and fix a source of
trivial, crash-class failures unrelated to clarification quality: the
Orchid-BCB and Orchid-BCB-Expand ambiguity injections rewrite BigCodeBench's
free-text instruction prompt, which never carried parameter type
information to begin with, causing the coding model to guess incorrect
types for under-specified parameters (e.g., assuming a millisecond-epoch
integer timestamp was a \texttt{datetime} object). Reinstating bare
parameter types, requirements, and a worked example -- without
reinstating the full parameter-description text that would re-leak the
injected ambiguity -- and holding the ranking/clarification pipeline
completely fixed, on the 106 problems where the exact round-1/round-2
clarified prompt text could be reconstructed without a new GPU run:
round-1 mean fractional pass rate rises from $0.796$ to $0.863$, with the
zero-pass (complete-crash) rate falling from $2.8\%$ to $0.0\%$; round-2
rises from $0.804$ to $0.854$, zero-pass from $3.8\%$ to $0.9\%$. This
improvement is concentrated almost entirely in Orchid-BCB-Expand and is
orthogonal to the ranking-quality question in
Section~\ref{sec:results:correctness}: it changes prompt content only, with
the identical clarifying question selected in both conditions.

\section{Discussion and Threats to Validity}
\label{sec:discussion}

\subsection{A Reversed Finding, Disclosed}

An earlier stage of this investigation, evaluated on a smaller
(121-group), more loosely filtered split, found the trained ranking probe
outperforming zero-shot self-rank by 13 points. That result did not
survive evaluation on the larger, more stringently filtered 1,229-group
pool reported in Section~\ref{sec:results:ranking}, where zero-shot wins
on 6 of 7 rows. We report both results and the reversal explicitly rather
than reporting only the final number, since a reversal under closer
scrutiny is itself informative: it indicates the earlier comparison was
underpowered or measuring an artifact of the looser filter, and a reader
relying on the earlier number alone would draw the wrong practical
conclusion about whether to deploy the probe in place of a zero-shot
baseline.

\subsection{Sample Size and Statistical Power}

Several of the correctness results in Section~\ref{sec:results:correctness}
are drawn from modest sample sizes -- most notably Orchid-BCB at
$n{=}21$ for the hit1-vs-pass-rate comparison. Its $+21.5$pp round-1 gap is
suggestive (Mann-Whitney $p = 0.057$) but not conclusive, and a swing of
one or two cases would move it substantially; we report the wide bootstrap
confidence interval alongside the point estimate for this reason rather
than treating a directionally-interesting result as a settled one. The
headline all-Orchid-combined result ($n{=}213$) is comparatively better
powered and reaches significance by two of three tests applied, but we
regard replication at larger $n$ as necessary before treating the specific
magnitude of the effect as established.

\subsection{Cross-Source Generalization}

The ranking probe's leave-one-source-out performance varies substantially
by held-out source (Table~\ref{tab:ranking}), transferring worst to
Orchid-BCB-Expand. We traced this to the injected ambiguous edit
constituting a small fraction of a long, free-form prompt for
BigCodeBench-derived sources (median edit fraction $2.1$--$2.3\%$ of
prompt words) compared to Orchid-HEval's shorter, more directly-edited
prompts (median $7.1\%$) -- a single pooled or last-token hidden-state
vector has proportionally less of its representational capacity available
to encode a small, diluted edit within a long prompt. This is consistent
with zero-shot self-rank, which re-reads the full prompt text at inference
time rather than compressing it into one fixed-size vector, not suffering
the same degradation on exactly this source. We consider this an
important limitation of single-vector hidden-state probing for long,
free-form specifications, rather than a limitation specific to our probe
architecture, and a natural target for future work (e.g., multi-token or
attention-weighted pooling that preserves more positional information about
where in a long prompt the disambiguating content lives).

\subsection{Choice of Probe-Host Model}

We use Qwen3-8B as the probe-host model throughout due to compute
constraints available to this study, not because it is necessarily the
best choice. A larger probe-host model might encode ambiguity and
candidate-quality information more richly, at the cost of increasing the
per-forward-pass cost that is central to this approach's efficiency
argument; conversely, a smaller model might preserve most of the cost
advantage while sacrificing some signal quality. We do not have data
characterizing this tradeoff and flag it as a direct avenue for future
work rather than assume our specific model choice generalizes.

\subsection{Metric Sensitivity}

Section~\ref{sec:results:correctness}'s central methodological finding --
that the fractional test-case pass rate recovers a significant
ranking-quality effect that the all-or-nothing pass metric does not --
depends on test suites that assign partial credit meaningfully, i.e.,
problems with multiple, at least partially independent test assertions.
For problems with very few test cases, the two metrics converge and this
distinction has little practical effect; we would expect the gap between
metrics to shrink correspondingly on benchmarks with coarser test suites.

\subsection{Absence of a Direct Repeated-Sampling Baseline}

We compare against zero-shot self-rank, which requires a single generation
call per problem, but we do not implement or benchmark a repeated-sampling,
self-consistency-based baseline in the style of ClarifyGPT's underlying
uncertainty-estimation mechanism~\cite{clarifygpt} directly against our
pipeline on the same problems. Any comparison to that class of method in
this paper is therefore a related-work comparison grounded in its published
methodology, not an empirical head-to-head measurement, and should be read
accordingly; we regard implementing and directly benchmarking such a
baseline as important future work to substantiate the cost-efficiency
argument beyond the zero-shot-vs-probe comparison reported here.

\subsection{Construct Validity of the Ambiguity Injections}

Both ClarifyCodeBench and the Orchid-derived sources construct ambiguity by
programmatically editing an originally unambiguous prompt, rather than
sourcing naturally occurring ambiguous specifications. This is a common and
practical approach to constructing matched clear/ambiguous pairs with known
ground truth, but it means our results characterize how well this pipeline
detects and resolves a specific, constructed style of ambiguity, and may
not transfer identically to naturally occurring specification gaps in
real development workflows.

\section{Conclusion and Future Work}
\label{sec:conclusion}

We investigated whether a lightweight linear probe over a small,
decoupled model's hidden states -- extracted without sampling or
generating any text -- can detect specification ambiguity in coding
prompts and rank candidate clarifying questions for an agentic
clarify-then-code pipeline, as a lower-cost alternative to uncertainty
estimation via repeated sampling of the coding model itself. Evaluated
against a zero-shot self-rank baseline under grouped, leave-one-source-out
splits across four benchmark sources, the probe reaches ranking quality
close to, and on its training-distribution source above, the baseline, at
substantially lower measured per-decision inference cost, framing this as
a cost-quality tradeoff rather than an unconditional accuracy improvement.
Deployed in a live, two-round elicit-rank-clarify-code loop against real
executable test suites, we found that ranking the labeled ambiguity
correctly is associated with a statistically significant round-1
correctness gain when measured by fractional test-case pass rate, though
not when measured by an all-or-nothing pass metric on the same data -- a
discrepancy we treat as evidence that correctness-metric choice materially
affects what an agentic clarification pipeline's evaluation can detect.
We further found that asking any clarifying question at all, independent
of ranking precision, is itself a robust and significant correctness
improvement, and that a cheap, targeted prompt augmentation addressing a
separate, orthogonal source of crash-class failures further improves
outcomes.

We reported several points where earlier-stage conclusions in this
investigation -- most notably an initial finding that the probe
outperformed zero-shot ranking -- did not survive evaluation at larger
scale and under stricter filtering, and disclose these reversals as part
of the record rather than reporting only the final numbers. We believe
this level of disclosure is necessary for agentic software engineering
research broadly, where pipelines with several interacting stages (elicit,
rank, clarify, generate, test) create many opportunities for an early,
underpowered, or confounded result to look like a finding.

Future work includes: benchmarking directly against a repeated-sampling,
self-consistency baseline on identical problems to substantiate the
cost-efficiency argument empirically rather than by citation; measuring
the actual batched $N$-candidate extraction cost against zero-shot
generation, rather than the bare-prompt-only forward pass reported here;
characterizing the tradeoff between probe-host model size and both ranking
quality and inference cost; and improving cross-source generalization for
long, free-form prompts where the ambiguous edit constitutes a small
fraction of the total prompt, which we identified as the ranking probe's
principal weakness.


\bibliographystyle{IEEEtran}
\bibliography{references}

\end{document}
